\documentclass[10pt,twocolumn]{article}

\usepackage[margin=0.72in,columnsep=0.25in]{geometry}
\usepackage[T1]{fontenc}
\usepackage[utf8]{inputenc}
\usepackage{lmodern}
\usepackage{microtype}
\usepackage{amsmath,amssymb}
\usepackage{booktabs}
\usepackage{array}
\usepackage{tabularx}
\usepackage{enumitem}
\usepackage{graphicx}
\usepackage{tikz}
\usetikzlibrary{arrows.meta,positioning,shapes.geometric,fit}
\usepackage{xcolor}
\usepackage{listings}
\usepackage{url}
\usepackage[hidelinks]{hyperref}
\usepackage[backend=biber,style=authoryear,natbib=true,maxcitenames=2,maxbibnames=99]{biblatex}
\addbibresource{references.bib}
\usepackage{caption}
\usepackage{fancyhdr}

\definecolor{capblue}{RGB}{24,78,119}
\definecolor{capgreen}{RGB}{42,125,79}
\definecolor{capgray}{RGB}{245,247,250}
\definecolor{capred}{RGB}{170,55,55}

\hypersetup{
  pdftitle={CAP: A Governance-First Wire Protocol for Distributed AI Agent Control Planes},
  pdfauthor={Yaron Torgeman},
  pdfsubject={Distributed AI agent governance and control-plane protocols},
  pdfkeywords={AI agents, governance, distributed systems, protocol, safety, authorization}
}

\lstset{
  basicstyle=\ttfamily\footnotesize,
  backgroundcolor=\color{capgray},
  frame=single,
  rulecolor=\color{black!15},
  breaklines=true,
  columns=fullflexible,
  keepspaces=true,
  showstringspaces=false,
  xleftmargin=2pt,
  xrightmargin=2pt
}

\pagestyle{fancy}
\fancyhf{}
\fancyhead[L]{CAP Technical Report}
\fancyhead[R]{June 2026}
\fancyfoot[C]{\thepage}
\renewcommand{\headrulewidth}{0.3pt}

\newcommand{\capname}{CAP}
\newcommand{\decision}{\mathcal{D}}
\newcommand{\jobs}{\mathcal{J}}
\newcommand{\workers}{\mathcal{W}}
\newcommand{\policies}{\mathcal{P}}
\newcommand{\states}{\mathcal{S}}

\title{\vspace{-1.4em}\textbf{CAP: A Governance-First Wire Protocol\\for Distributed AI Agent Control Planes}}
\author{Yaron Torgeman\\\small Cordum\\\small \texttt{admin@cordum.io}}
\date{Technical Report, June 22, 2026\\\small Artifact version: CAP commit \texttt{df2af475}; Cordum commit \texttt{6232e9d}}

\begin{document}
\raggedbottom
\maketitle
\thispagestyle{fancy}

\begin{abstract}
Autonomous agents increasingly execute work across distributed worker fleets, yet most agent protocols standardize tool invocation or peer messaging rather than the control plane that admits, governs, schedules, and reconciles consequential work. We present the Cordum Agent Protocol (CAP), a governance-first wire protocol for distributed AI agent control planes. CAP represents each unit of work as a versioned job contract containing identity, capability metadata, pointer-based context, workflow lineage, and resource budgets. A scheduler evaluates policy before dispatch, selects eligible workers from expiring capacity advertisements, and advances an explicit lifecycle over transports that may redeliver or reorder messages. We formalize invariants for mediated dispatch, capacity-bounded placement, monotone state transitions, idempotent redelivery, and control/data separation. We describe Cordum, the reference implementation over NATS and Redis, and audit CAP's public artifact history from v0.1.0 on December 11, 2025 through the current append-only wire version. CAP predates the 2026 cluster of execution-governance specifications as a public, versioned protocol for distributed agent fleets. We do not claim invention of reference monitoring, authorization, contracts, or audit logging; CAP's contribution is their integration into an interoperable distributed job lifecycle. The specification, SDKs, fixtures, and reference implementation are openly available.
\end{abstract}

\section{Introduction}

Agent systems have moved from producing text to producing effects. A single user request can fan out across model workers, tool runners, cloud APIs, databases, and sub-agents. The resulting system is not merely an inference service; it is a distributed execution environment whose failures include duplicate actions, stale placement decisions, cross-tenant confusion, unbounded resource use, policy bypass, partial execution, and irreconcilable state.

Tool protocols such as the Model Context Protocol (MCP) standardize how models discover and invoke tools, while peer protocols such as Agent2Agent (A2A) standardize inter-agent communication \citep{mcp2024,a2a2025}. These are valuable interoperability layers, but neither is a complete control-plane contract. A production control plane must answer a different set of questions: What is the unit of schedulable work? Which principal and tenant caused it? What budget applies? Which policy snapshot authorized dispatch? Which workers are currently eligible? How does the system recover from duplicate delivery or a worker disappearing after dispatch? Where do large or sensitive payloads reside? How is a workflow reconstructed without inspecting model reasoning?

CAP was created to make those concerns part of the wire contract rather than framework-specific application logic. The repository history records v0.1.0 on December 11, 2025 with a versioned envelope, jobs and results, worker heartbeats, a pre-dispatch Safety Kernel, human escalation, lifecycle states, pointer-based payload separation, and workflow lineage \citep{cap_v010}. It records a release the next day adding first-class token and deadline budgets together with tenant and principal identity \citep{cap_v200}. Because repository tags and author dates remain maintainer-controlled, we do not use those dates alone as proof of public availability. An independent Go package-registry record establishes CAP distribution by January 9, 2026 \citep{cap_pkg_registry}. Cordum is the production-oriented reference implementation \citep{cordum_current,cordum_disclosure2026}.

The design inherits established security principles. The Safety Kernel is a reference-monitor pattern, not a newly invented security primitive \citep{anderson1972reference,saltzer1975protection}. Job delegation and distributed coordination have deep roots, including the Contract Net Protocol \citep{smith1980contractnet}. CAP's contribution is narrower and operational: it binds governance, capacity, lifecycle, and delivery semantics into one interoperable protocol for heterogeneous agent fleets.

This paper makes four contributions:
\begin{enumerate}[leftmargin=*,nosep]
  \item \textbf{A control-plane wire model.} CAP defines a transport-neutral envelope and job lifecycle for gateways, schedulers, workers, orchestrators, and policy services.
  \item \textbf{Governance-bearing job contracts.} Jobs carry identity, capability, risk, budget, pointer, and workflow metadata, while policy decisions can deny, defer, throttle, redact, quarantine, or tighten execution constraints.
  \item \textbf{Failure-aware distributed semantics.} Expiring heartbeats, explicit states, idempotency, retries, cancellation, and at-least-once delivery assumptions are part of the specification.
  \item \textbf{A public artifact and reference implementation.} CAP includes normative documents, protocol buffers, conformance levels, deterministic fixtures, multi-language SDKs, and Cordum as a code-backed implementation.
\end{enumerate}

\section{Scope and Threat Model}

\subsection{System model}

CAP assumes five logical roles:
\begin{itemize}[leftmargin=*,nosep]
  \item a \emph{gateway} authenticates callers, stores input payloads, and submits jobs;
  \item a \emph{scheduler} owns lifecycle state, invokes policy, and dispatches work;
  \item a \emph{Safety Kernel} returns deterministic policy decisions and constraints;
  \item a \emph{worker} advertises capacity, executes a job, stores output, and emits a result;
  \item an \emph{orchestrator} creates related jobs and aggregates workflow state.
\end{itemize}
The roles may be separate processes, replicated services, or embedded components. CAP standardizes their messages and minimum behavior rather than prescribing a deployment topology.

\subsection{Adversary and failure assumptions}

We consider an agent or upstream model that may generate unsafe, malformed, excessive, or adversarial requests. We also consider ordinary distributed-system faults: duplicate and reordered delivery, stale heartbeats, worker crashes, scheduler restarts, policy-service timeouts, partial persistence, and network partitions. A caller may attempt tenant confusion, replay, or capability escalation.

The trusted computing base includes the enforcement path that prevents a worker from receiving governed work except through the scheduler, the policy service used for a decision, and the authenticated transport and state store. CAP cannot make enforcement non-bypassable if a worker exposes an unaudited side channel. Likewise, a declared sandbox or tool constraint is only as strong as the runtime that enforces it.

\subsection{Non-goals}

CAP does not specify model alignment, prompt formats, reasoning traces, a universal policy language, or a global identity provider. It does not promise exactly-once execution over an asynchronous network. The current wire includes packet signatures and runtime tokens, but CAP does not yet define a complete cross-organization trust root or a Merkle-based evidence ledger. Result-content governance is implemented by Cordum as an extension, but it is not yet a required CAP conformance feature.

\input{sections/03-protocol.tex}
\section{Governance and Distributed-System Invariants}

We separate \emph{protocol invariants} from \emph{deployment assumptions}. A conforming message format alone cannot prevent bypass; the scheduler and executors must make the protocol path authoritative.

\subsection{Mediated dispatch}

Let $P(J,\pi,s)$ evaluate job $J$ under policy snapshot $\pi$ and observable state $s$. Let $d$ be the resulting decision and $K$ its constraints. Dispatch is safe only if
\begin{equation}
\begin{aligned}
Dispatch(J,w) \Rightarrow d \in \{&ALLOW,\\
&ALLOW\_WITH\_CONSTRAINTS\}.
\end{aligned}
\end{equation}
and $w \in E(J,t)$. A human decision resolves a deferred request; it is not equivalent to a generic credential that authorizes unrelated future jobs.

Cordum strengthens this property by binding approval to the stored job representation and policy snapshot before requeueing. CAP carries the snapshot and approval reference needed for equivalent implementations, although the exact hash-binding procedure is not yet normative.

\subsection{Constraint monotonicity}

When both a requester and policy specify upper bounds, an enforcement implementation should derive an effective bound no weaker than either source. For a scalar limit $b$,
\begin{equation}
  b_{effective} = \min(b_{request}, b_{policy})
\end{equation}
when both are present. Cordum applies this tightening to runtime deadlines and enforces policy-provided retry and concurrency caps. CAP currently specifies the fields but leaves complete multi-dimensional merge semantics to a future normative profile.

\subsection{Monotone lifecycle}

Let $\prec$ be the permitted state relation. For every accepted event $e_i$,
\begin{equation}
  S_{i+1} = S_i \quad \text{or} \quad S_i \prec S_{i+1}.
\end{equation}
Duplicate messages may leave the state unchanged. A result for an already terminal job cannot move it back to a non-terminal state. This property makes redelivery safe and supports deterministic reconciliation after restarts.

\subsection{Idempotent redelivery}

For a fixed idempotency key $k$, repeated submission must not create unbounded side effects. CAP requires consumers to tolerate redelivery but does not claim network-level exactly once. Implementations use locks, compare-and-set state, or result caching to approach effect-once behavior for idempotent workloads.

\subsection{Capacity-bounded dispatch}

A scheduler must not deliberately dispatch to a worker whose latest valid advertisement reports no available slot or lacks required capabilities. Since heartbeats can be delayed, this is a placement safety property over the scheduler's latest authenticated view, not a guarantee about instantaneous physical load.

\subsection{Control/data separation}

For any bus packet $p$, large input and output payloads are referenced rather than embedded:
\begin{equation}
  payload(p) = metadata \cup pointers,
\end{equation}
subject to small inline fields defined by a profile. This supports independent retention, redaction, and storage authorization policies.

\section{Cordum Reference Implementation}

Cordum realizes CAP as a set of Go services over NATS and Redis \citep{cordum_current}. The API gateway authenticates callers, enforces tenant context, writes input payloads, and publishes job requests. The scheduler persists state, calls the Safety Kernel, resolves pool and worker eligibility, dispatches, and reconciles stale work. The Safety Kernel evaluates signed and versioned policies and returns constraints, approvals, and remediation options. The workflow engine expands DAG steps into CAP jobs and consumes job results. External workers use CAP SDKs and remain outside the control-plane repository.

\subsection{Failure handling}

With JetStream enabled, durable subjects use explicit acknowledgment and negative acknowledgment with delay. The scheduler protects critical sections with per-job Redis locks and renews lock leases during long policy or routing decisions. A pending replayer recovers submissions that did not reach dispatch, while a reconciler marks stale dispatched or running jobs as timed out. Retryable failures avoid immediate dead-letter handling; fatal failures can trigger reverse-order compensation.

The worker registry is intentionally ephemeral. Heartbeats populate an in-memory map with a default 30-second time-to-live. This removes a persistent worker database from the dispatch path; a replicated snapshot can accelerate warm start without becoming the authority for liveness.

\subsection{Policy and output extensions}

Cordum persists decision records containing policy snapshots, rule identifiers, reasons, constraints, and approvals. It also implements a post-execution output-policy layer with allow, redact, deny, and quarantine decisions. That layer is useful evidence that CAP can support multiple governance boundaries, but the present paper does not claim it as part of CAP's minimum wire conformance.

\subsection{Security boundary}

The reference implementation enforces API authentication, tenant checks, TLS requirements in production, policy signature verification, URL allowlists for remote policy loading, bounded I/O, and explicit timeouts. These are implementation properties, not automatic consequences of speaking CAP. A minimal CAP implementation can be interoperable while still being insecure; conformance and hardening are distinct dimensions.

\section{Artifact Evaluation}

This report evaluates the protocol as a public systems artifact rather than presenting unverified throughput claims. All observations are tied to immutable tags or pinned commits.

\subsection{Historical continuity}

Table~\ref{tab:timeline} separates maintainer-controlled repository history from independently recorded publication evidence. Git commits and tags are useful provenance, but their dates and ref targets are not immutable to repository administrators. The earliest independent anchor currently cited by this report is the Go package registry, which records CAP v2.0.8 as published on January 9, 2026 \citep{cap_pkg_registry}. That date predates Agent Contracts, Faramesh, AARM, OpenPort, OAP, SARC, and C-Trace, but it does not independently establish that the December 11 tag was publicly reachable on that day. We therefore describe December as repository history and January 9 as the externally anchored public date until stronger evidence is archived.

\begin{table}[t]
\centering
\caption{Selected CAP milestones and evidence classes.}
\label{tab:timeline}
\small
\begin{tabular}{@{}p{0.22\columnwidth}p{0.69\columnwidth}@{}}
\toprule
Date & Evidence and milestone \\
\midrule
2025-12-11 & Repository history: v0.1.0 records envelope, jobs/results, heartbeats, pools, pre-dispatch Safety Kernel, human deferral, lifecycle, pointers, and workflows. \\
2025-12-12 & Repository history: v2.0.0 records token/deadline budgets, memory hints, tenant and principal identity, and labels. \\
2026-01-03 & Repository history: structured capability/risk metadata, policy snapshots, constraints, simulation/explanation surfaces, and concurrency caps. \\
2026-01-09 & Independent registry: pkg.go.dev records CAP v2.0.8 as published. \\
2026-01-23 & Repository disclosure: Cordum CAP v2 design and implementation mapping. \\
2026-06-21 & Pinned current CAP and Cordum commits used by this report. \\
\bottomrule
\end{tabular}
\end{table}

\subsection{Priority-evidence policy}

A stronger historical claim should additionally cite a GitHub organization audit-log visibility event, an OpenTimestamps proof, a Software Heritage identifier, and/or an Internet Archive capture. Those artifacts are intentionally not represented as complete in this version. The manuscript's claim policy is therefore: (1) use December 2025 only for repository-history statements; (2) use January 9, 2026 as the earliest independently anchored public distribution date currently in evidence; and (3) update the claim only after immutable third-party anchors are archived with the artifact bundle.

\subsection{Compatibility surface}

The current specification contains 17 numbered normative documents, uses RFC 2119 requirement language \citep{rfc2119}, and separates wire version from SDK release version. The wire remains version 1 because evolution has been append-only. Deterministic binary fixtures are consumed by SDK tests to detect serialization drift. CAP also defines CORE, STANDARD, and FULL conformance levels so an implementation can state which lifecycle, security, and orchestration capabilities it supports.

\subsection{Implementation traceability}

Table~\ref{tab:traceability} maps major properties to public artifacts. This is not a claim that every deployment is secure; it shows that the protocol concepts have concrete code paths and tests in the reference implementation.

\begin{table*}[t]
\centering
\caption{Traceability from protocol property to CAP and Cordum artifacts.}
\label{tab:traceability}
\small
\begin{tabular}{@{}p{0.18\textwidth}p{0.32\textwidth}p{0.40\textwidth}@{}}
\toprule
Property & CAP artifact & Cordum realization \\
\midrule
Pre-dispatch governance & \texttt{safety.proto}; normative Safety specification & Scheduler Safety client; Safety Kernel; decision persistence; approval handling \\
Budgeted job contract & \texttt{job.proto} \texttt{Budget}; \texttt{JobMetadata} & Deadline, retry, artifact, and tenant concurrency enforcement \\
Capacity-aware routing & \texttt{heartbeat.proto}; pool semantics & TTL worker registry; least-loaded and capability-aware strategy \\
At-least-once safety & Transport profile; idempotency requirements & JetStream ack/nak; per-job locks; pending replay; timeout reconciliation \\
Pointer separation & Memory pointer specification & Redis context, result, and artifact stores; gateway pointer APIs \\
Workflow lineage & Parent/workflow/step fields; compensation & DAG workflow engine; retries; approval steps; saga rollback \\
Cross-SDK compatibility & Stable envelope; append-only versioning & Deterministic fixtures and generated SDKs \\
\bottomrule
\end{tabular}
\end{table*}

\subsection{What is deliberately not measured}

The repository contains historical benchmark documents, but this report does not import those values because the manuscript artifact does not yet bundle the raw traces, pinned load generator, complete hardware image, warm-up procedure, and confidence intervals required for independent reproduction. A future empirical revision should report at least: Safety Kernel decision latency (p50/p95/p99 over a declared request corpus), scheduler dispatch throughput, duplicate-delivery convergence, stale-heartbeat exclusion, and a worked destructive-action denial with its decision and audit records. These experiments must use a tagged harness and publish raw results. Excluding unsupported numbers is a deliberate validity choice, not an absence of implementation.

\section{Related Work}

\subsection{Foundational control and coordination}

Reference monitoring, complete mediation, least privilege, and fail-safe defaults long predate agent systems \citep{anderson1972reference,saltzer1975protection}. Contract-based multi-agent coordination dates at least to the Contract Net Protocol \citep{smith1980contractnet}. CAP applies these ideas to a modern, heterogeneous execution fabric; it does not rename them as new inventions.

\subsection{Agent communication protocols}

MCP standardizes model-to-tool context, resources, prompts, and tool calls \citep{mcp2024}. A2A standardizes discovery and task-oriented communication between agents \citep{a2a2025}. CAP is complementary: it standardizes the governance-bearing distributed job lifecycle between gateways, schedulers, policy services, and worker pools. MCP may run inside a CAP worker, and an A2A endpoint may be exposed as a CAP-managed capability.

\subsection{Governance frameworks and bounded execution}

AGENTSAFE provides a broad lifecycle governance and assurance framework and predates CAP's repository-recorded December release \citep{khan2025agentsafe}. Agent Contracts formalizes resource-bounded execution and delegated budget conservation \citep{ye2026agentcontracts}. Policies on Paths makes the execution path, partial history, proposed action, identity, and organizational state the central object of runtime governance \citep{kaptein2026policies}. CAP contributes an interoperable wire representation and reference implementation; it does not claim Agent Contracts' formal delegation laws or the path-level risk formalization of Policies on Paths.

\subsection{Pre-action and execution authorization}

Faramesh defines an Action Authorization Boundary and Canonical Action Representation \citep{fatmi2026faramesh}. PCAS compiles Datalog-derived policies into an instrumented reference monitor and reports policy compliance increasing from 48\% to 93\% with zero policy violations in instrumented customer-service runs \citep{palumbo2026pcas}. AARM specifies interception, contextual policy evaluation, enforcement, and receipts across several deployment architectures \citep{errico2026aarm}. OpenPort specifies a governed tool gateway with authorization-dependent discovery, human-reviewed writes, and state witnesses \citep{zhu2026openport}. OAP synchronously intercepts tool calls and emits signed authorization records \citep{uchibeke2026oap}. These works provide richer action-level authorization models or measurements than CAP's initial protocol. CAP differs in scope: its principal abstraction is a distributed fleet job with routing, liveness, budget, lifecycle, and workflow semantics.

\subsection{Multi-stage governance and verifiable history}

SARC compiles constraints into pre-action, action-time, post-action, and escalation sites \citep{besanson2026sarc}. C-Trace monitors tool invocations and model outputs for runtime compliance \citep{kahani2026ctrace}. Right to History implements a local sovereignty kernel with Merkle-verifiable logs \citep{zhang2026righthistory}, while OpenKedge cryptographically links intent, context, policy, execution bounds, and outcomes \citep{he2026openkedge}. CAP's initial artifact does not provide equivalent evidence-chain guarantees. Its distinctive contribution is the open distributed control-plane wire contract on which such mechanisms can be carried or composed.

\begin{table*}[t]
\centering
\caption{Qualitative scope comparison. A check mark means the item is a primary, explicit contribution of the cited work, not merely an implementable extension.}
\label{tab:related}
\scriptsize
\setlength{\tabcolsep}{3pt}
\begin{tabularx}{\textwidth}{@{}l>{\centering\arraybackslash}X>{\centering\arraybackslash}X>{\centering\arraybackslash}X>{\centering\arraybackslash}X>{\centering\arraybackslash}X@{}}
\toprule
Work & Pre-execution policy & Fleet scheduling / liveness & Budget in work contract & Multi-stage output control & Cryptographic evidence focus \\
\midrule
CAP & $\checkmark$ & $\checkmark$ & $\checkmark$ & extension & optional packet signing \\
Agent Contracts & contract semantics & -- & $\checkmark$ & -- & -- \\
Policies on Paths & path policy & fleet framing & organizational risk & path-dependent & -- \\
PCAS & $\checkmark$ & causal dependency graph & -- & information-flow paths & enforcement traces \\
Faramesh & $\checkmark$ & -- & -- & -- & provenance artifacts \\
OpenPort & $\checkmark$ & -- & scoped limits & -- & audit/state witness \\
AARM & $\checkmark$ & -- & contextual & receipts & tamper-evident receipts \\
OAP & $\checkmark$ & -- & policy limits & -- & signed audit record \\
SARC & $\checkmark$ & workflow propagation & constraints & $\checkmark$ & trace correspondence \\
Right to History & human approval & local kernel & energy budget & -- & $\checkmark$ \\
OpenKedge & $\checkmark$ & conflict arbitration & execution bounds & outcomes & $\checkmark$ \\
C-Trace & $\checkmark$ & -- & -- & $\checkmark$ & compliance traces \\
\bottomrule
\end{tabularx}
\end{table*}

\section{Limitations and Research Agenda}

First, CAP standardizes messages and expected behavior, but non-bypassability depends on deployment architecture. Direct worker credentials or side channels can undermine the control plane. Second, sandbox, toolchain, and network constraints require enforcement by the selected runtime; a protobuf field alone is not isolation. Third, CAP provides at-least-once tolerance rather than exactly-once execution. Non-idempotent effects require downstream idempotency keys, transactional APIs, or compensating actions.

Fourth, current packet signatures and tokens do not constitute a complete multi-tenant evidence chain. Integrating DSSE-style envelopes, transparent checkpoints, or customer-controlled witnesses is future work. Fifth, output governance is not yet a required CAP profile. A staged-release state machine could make post-execution inspection interoperable. Sixth, the formal semantics in this report should be encoded in TLA+, PlusCal, or another model checker and validated under crash, partition, and redelivery schedules.

The immediate empirical agenda is a tagged, one-command harness with raw data for: (1) policy and dispatch latency distributions; (2) scaling with workers, pools, and tenants; (3) invariant preservation under duplicate and reordered delivery; (4) stale-heartbeat and policy-outage behavior; and (5) release safety for output redaction and quarantine.

\section{Conclusion}

CAP treats an autonomous agent fleet as a governed distributed system rather than a collection of model tool calls. Its wire contract connects identity, budgets, policy, live capacity, pointers, lifecycle, and workflows while remaining independent of model and transport. The repository history records this integrated control-plane design in December 2025, and an independent package-registry record establishes public distribution by January 9, 2026. The former is maintainer-controlled provenance; the latter is the current external anchor. Both are narrower claims than invention of the underlying security and coordination concepts, which are established prior art. Cordum shows that the protocol maps to a production-oriented implementation with explicit failure handling and extensible policy. The next step is not broader firstness language; it is stronger formalization, reproducible fault evaluation, and interoperable evidence and output-release profiles.

\section*{Availability and Ethics Statement}

CAP is available under Apache-2.0 at \url{https://github.com/cordum-io/cap}. Cordum is available at \url{https://github.com/cordum-io/cordum}. The paper source pins exact commits in its artifact manifest. This report uses no human-subject data, customer data, or proprietary evaluation corpus. It makes no claim of peer review and should be cited as a technical report or arXiv preprint until accepted by a venue.

\section*{Acknowledgments}

The author thanks the CAP and Cordum contributors whose implementation, testing, documentation, and protocol reviews made the public artifact possible.


\printbibliography

\end{document}
